\DocumentMetadata{
  pdfversion=2.0,
  pdfstandard=ua-2,
  lang=en-US,
  tagging=on,
  tagging-setup={math/setup=mathml-SE,tabsorder=structure}
}
\documentclass[11pt,letterpaper]{article}
\usepackage[margin=0.68in,includefoot]{geometry}
\usepackage{newcomputermodern}
\usepackage{amsmath,amscd,mathtools}
\providecommand{\square}{\mdlgwhtsquare}
\usepackage[hidelinks]{hyperref}
\hypersetup{
  pdftitle={Algebra PhD Exam, September 1995},
  pdfauthor={University of Florida Department of Mathematics},
  pdfsubject={Graduate mathematics examination},
  pdfdisplaydoctitle=true,
  bookmarksopen=true
}
\setcounter{secnumdepth}{0}
\setlength{\parindent}{0pt}
\setlength{\parskip}{0.28em}
\setlength{\emergencystretch}{3em}
\setlength{\leftmargini}{2.15em}
\setlength{\leftmarginii}{2.65em}
\setlength{\leftmarginiii}{3em}
\renewcommand{\labelenumi}{\textbf{\theenumi.}}
\pagestyle{empty}
\raggedbottom
\allowdisplaybreaks
\newenvironment{examproblems}[1]{%
  \begin{enumerate}%
  \setcounter{enumi}{#1}%
  \setlength{\topsep}{0.35em}%
  \setlength{\partopsep}{0pt}%
  \setlength{\itemsep}{0.48em}%
  \setlength{\parsep}{0.18em}%
}{\end{enumerate}}
\newenvironment{examsubparts}{%
  \begin{enumerate}%
  \setlength{\topsep}{0.18em}%
  \setlength{\partopsep}{0pt}%
  \setlength{\itemsep}{0.16em}%
  \setlength{\parsep}{0.1em}%
}{\end{enumerate}}
\newenvironment{examinstructions}{%
  \par\begingroup\itshape
}{%
  \par\endgroup\vspace{0.18em}
}
\makeatletter
\renewcommand\section{\@startsection{section}{1}{0pt}%
  {-1.25ex plus -.25ex minus -.1ex}%
  {0.55ex plus .1ex}%
  {\normalfont\large\bfseries}}
\renewcommand{\@maketitle}{%
  \newpage\null\vskip -1.2em%
  {\footnotesize\bfseries
    \href{https://www.ufl.edu/}{University of Florida}, \href{https://math.ufl.edu/}{Department of Mathematics}\par}%
  \vskip 0.18em%
  \tagstructbegin{tag=Title}%
    {\LARGE\bfseries\href{https://gma.math.ufl.edu/exams/algebra-phd/}{Algebra PhD Exam}, September 1995\par}%
  \tagstructend%
  \vskip 0.35em%
  \tagmcbegin{artifact}%
    \hrule height 1pt%
  \tagmcend%
  \vskip 0.55em%
}
\makeatother
\title{Algebra PhD Exam, September 1995}
\author{}
\date{}
\begin{document}
\maketitle
\thispagestyle{empty}
\begin{examinstructions}
Answer seven problems. Write your answers clearly in complete English sentences. You may quote results (within reason) as long as you state them clearly.
\end{examinstructions}
\begin{examproblems}{0}
\item
This problem has two parts.
\begin{examsubparts}
\item[\textnormal{(a)}]
Give the definitions of a solvable group and a nilpotent group.
\item[\textnormal{(b)}]
Let~\(G\) be a finite solvable group. Let \(F(G)\) be the largest normal nilpotent subgroup of~\(G\). Prove that if~\(G\) is not abelian, then \(F(G)>Z(G)\), where \(Z(G)\) is the center of~\(G\) (ie., \(F(G)\) is bigger but not equal to \(Z(G)\).
\end{examsubparts}
\item
A proper subgroup~\(M\) of a group~\(G\) is maximal if whenever \(M<H\leq G\), we have \(H=M\) or \(H=G\). Suppose that~\(G\) is a finite group has only one maximal subgroup. Prove that~\(G\) is cyclic of prime power order.
\item
Let~\(A\) be a commutative Artinian ring with identity. Prove the following statements.
\begin{examsubparts}
\item[\textnormal{(a)}]
All prime ideals are maximal.
\item[\textnormal{(b)}]
There is only finite number of prime, or maximal ideals.
\item[\textnormal{(c)}]
The ideal~\(N\) of nilpotent elements in~\(A\) is a nilpotent ideal.
\end{examsubparts}
\item
A module~\(F\) over a commutative ring with identity is called flat if for any short exact sequence \(0\longrightarrow E'\longrightarrow E\longrightarrow E''\longrightarrow 0\), the sequence \(0\longrightarrow F\otimes E'\longrightarrow F\otimes E\longrightarrow F\otimes E''\longrightarrow 0\) is exact.
\begin{examsubparts}
\item[\textnormal{(a)}]
Give an example of a non-flat module.
\item[\textnormal{(b)}]
Let \(0\longrightarrow F'\longrightarrow F\longrightarrow F''\longrightarrow 0\) be an exact sequence of modules over a commutative ring with identity. Suppose \(F''\) is a flat module. Prove that \(F'\) is flat if and only if~\(F\) is flat.
\end{examsubparts}
\item
This problem has two parts.
\begin{examsubparts}
\item[\textnormal{(a)}]
State the Jacobson Density Theorem.
\item[\textnormal{(b)}]
Using Jacobson’s density Theorem to prove the following Wedderburn’s Theorem: Let~\(R\) be a ring, and~\(E\) a simple, faithful module over~\(R\). Let \(D=\operatorname{End}_R(E)\), and assume that~\(E\) is a finite dimension over~\(D\). Then \(R=\operatorname{End}_D(E)\).
\end{examsubparts}
\item
Let~\(A\) be a commutative ring with identity, let~\(M\) be an~\(A\)-module, and let~\(N\) be a submodule of~\(M\).
\begin{examsubparts}
\item[\textnormal{(a)}]
Define what it means for~\(M\) to be Noetherian.
\item[\textnormal{(b)}]
Prove that if~\(M\) is Noetherian then the quotient \(M/N\) must also be Noetherian.
\item[\textnormal{(c)}]
Prove that a submodule of a Noetherian module is Noetherian.
\end{examsubparts}
\item
Let~\(A\) be a commutative ring with identity. Let~\(M\) be a~\(A\)-module.
\begin{examsubparts}
\item[\textnormal{(a)}]
Give the definition of an associated prime of~\(M\).
\item[\textnormal{(b)}]
Prove that if~\(A\) is Noetherian and \(M\neq 0\), then~\(M\) has an associated prime.
\end{examsubparts}
\item
Let~\(p\) be a prime and let \(U_p\) be the subgroup of the group of complex roots of unity consisting of those roots of unity whose order is a power of~\(p\).
\begin{examsubparts}
\item[\textnormal{(a)}]
Show that the map \(x\longmapsto x^n\) is a surjective endomorphism of \(U_p\), which is an automorphism if and only if~\(p\) does not divide~\(n\).
\item[\textnormal{(b)}]
Let~\(L\) denote the localization of the integers at the prime~\(p\). Deduce from (a) that \(U_p\) is a~\(L\)-module for which \((p)U_p=U_p\).
\item[\textnormal{(c)}]
State Nakayama’s Lemma and explain why the module \(U_p\) is not a counterexample.
\end{examsubparts}
\item
Consider the polynomial over the field \(F_2\) with two elements: 
\[
g(x)=x^4+x^3+x^2+x+1
\]
\begin{examsubparts}
\item[\textnormal{(a)}]
Prove that \(g(x)\) is irreducible over \(F_2\).
\item[\textnormal{(b)}]
Let~\(K\) be a splitting field for \(g(x)\) over \(F_2\) and let \(r\in K\) be a root of \(g(x)\). Factor \(g(x)\) into irreducibles over \(F_2(r)\).
\item[\textnormal{(c)}]
Show that \(K=F_2(r)\).
\item[\textnormal{(d)}]
Find the Galois group \(\operatorname{Gal}(K/F_2)\).
\end{examsubparts}
\item
Let \(K/k\) be a cyclic extension of fields with finite Galois group \(\langle\sigma\rangle\).
\begin{examsubparts}
\item[\textnormal{(a)}]
State the Hilbert 90 theorem (the multiplicative or the additive form).
\item[\textnormal{(b)}]
Suppose \([K:k]=n\) is prime to the \(\operatorname{Char}(k)\) and~\(k\) contains a primitive~\(n\)-th root of unity. Prove that there exists \(\alpha\in K\) such that \(K=k(\alpha)\) and \(\operatorname{Irr}(\alpha,k,x)\) (i.e., the irreducible polynomial of~\(\alpha\) over~\(k\)) \(=x^n-a\) for some \(a\in k\).
\end{examsubparts}
\item
Let \(G=\langle g\rangle\) be a finite cyclic group and regard~\(\mathbb{Z}\) (the integers) as a trivial module of~\(G\) (i.e., \(\mathbb{Z}[G]\)-module). Compute for \(n\geq 1\), \(H^n(G,\mathbb{Z})\). Justify your answer.
\end{examproblems}
\end{document}
